
The rise of foundation models changed not just what AI systems can accomplish but how researchers measure progress. This study identifies two statistically significant structural breaks in benchmark concentration---April 2019 and March 2021---coinciding with the emergence of transformative models such as GPT-2, BERT, GPT-3, and ViT. These breaks mark a phase transition in how the research community evaluates AI systems.

Understanding benchmark dynamics matters because benchmark choice shapes what problems researchers pursue. When 26.5\% of researcher attention flows toward emergent-capability benchmarks such as MMLU, BIG-Bench, and HumanEval---up from 7.5\% before 2021---the entire field's trajectory shifts. Without understanding these dynamics, it becomes difficult to distinguish genuine AI progress from benchmark artifact.

\subsubsection{The Problem of Static Analysis}

Prior work has documented benchmark concentration. Koch et al. (2021) found Gini coefficients of 0.6--0.7 for 2015--2020, demonstrating that research concentrates on fewer datasets over time. However, their analysis provides static snapshots ending at 2020---precisely when foundation models began reshaping the landscape.

The deeper problem is that foundation models may have coincided with a structural break in benchmark dynamics, fundamentally altering how concentration evolves. Detecting structural breaks requires time-series methods (change-point detection, trend analysis) rather than periodic snapshots. This methodological gap has left a critical question unanswered: did foundation models coincide with a measurable phase transition in benchmark usage patterns?

\subsubsection{Contributions}

\begin{enumerate}
\item \textbf{First quantitative evidence of phase transition in benchmark dynamics.} Using PELT change-point detection, two structural breaks (April 2019, March 2021) were identified with BIC improvement of 17.13 over a monotonic trend, providing statistical evidence that foundation model emergence coincided with measurable ecosystem restructuring.
\end{enumerate}

\begin{enumerate}
\item \textbf{A complete mechanism verification chain.} Five causal steps from foundation model emergence through benchmark ecosystem restructuring were tested, validating four of five mechanism hypotheses. The fifth (modality divergence) was refuted, yielding the unexpected finding that modality dynamics were always independent.
\end{enumerate}

\begin{enumerate}
\item \textbf{Quantified attention shift from traditional to emergent benchmarks.} Emergent-capability benchmark share increased from 7.53\% to 26.54\% post-2021 ($\chi^2$ = 1025.23, p < $10^{-224})$, while traditional benchmarks persisted with 47,068 papers but reduced relative dominance (11.70\% share).
\end{enumerate}
